\documentclass[class=scrbook, crop=false]{standalone}
\usepackage[subpreambles=true]{standalone}
\ifstandalone
    \input{../settings+/settings}
\fi

% ----------------------------------------------------------------------------
%                               Implementation
% ----------------------------------------------------------------------------
\begin{document}

\chapter{A Benchmark for Indefinite OCP Quadratic Programs} % Outline text
In order to evaluate the performance of the proposed Riccati-based primal regularization, 
especially in the context of the subproblems of SQP in NMPC.
In this chapter, we aim at constructing a benchmark of indefinite OCP QPs that arise in the SQP iterations of NMPC problems.
We will first introduce the construction of the QP benchmark, 
including the formulation of QPs in benchmark, the OCP test problems and the strategy for collecting indefinite QPs in Section \ref{sec:benchmark_construction}.
Then we will analyze the collected problem instances to gain a better understanding of characteristics
regarding definiteness in the indefinite OCP QP benchmark in Section \ref{sec:benchmark_analysis}.

\section{Benchmark Construction}
\label{sec:benchmark_construction}
In this section, we will introduce the construction of the indefinite QP benchmark.
We will first introduce the formulation of QPs in benchmark,
which originate from the SQP iterations of the OCPs, and then we will introduce the eight OCP test problems used to construct the benchmark,
which are formulated as fixed-time regulator problems with quadratic stage costs and box constraints on the controls.
Then the collection strategy is designed to ensure that the collected QPs are representative of the indefinite QPs that arise in practice, 
and that they share the mathematical structure of the OCPs.


\subsection{Formulation of QP subproblems from OCPs}

\label{subsec:formulation}
We will start from the formulation of the OCPs to obtain the resulting QP subproblems that arise in the SQP iterations.
In the benchmark, all optimal control problems share the same structure: 
a fixed-horizon, direct multiple-shooting \cite{bock1984multiple} optimal control problem 
with a quadratic stage cost that penalizes deviations from given state and control references.
The mathematical formulation of the OCPs is given as follows:
\begin{equation}
\label{eq:ocp}
\begin{aligned}
\min_{\substack{x_0,\dots,x_N\\ u_0,\dots,u_{N-1}}}\ \
& \sum_{k=0}^{N-1}\frac{1}{2}
  \begin{bmatrix} x_k-x_k^{\mathrm{ref}} \\ u_k-u_k^{\mathrm{ref}} \end{bmatrix}^{\!\top}
  \begin{bmatrix} \mathcal{Q} & \mathcal{S}_k^{\top} \\ \mathcal{S}_k & \mathcal{R} \end{bmatrix}
  \begin{bmatrix} x_k-x_k^{\mathrm{ref}} \\ u_k-u_k^{\mathrm{ref}} \end{bmatrix}
  +\tfrac12\|x_N-x_N^{\mathrm{ref}}\|_{\mathcal{Q}_N}^2\\
\text{s.t.}\ \
& r(x_0,x_N) = 0,\\
& x_{k+1}=f_k(x_k,u_k),  \quad  k=0,\dots,N-1,\\
& \underline{g} \le \mathcal{C}_k x_k + \mathcal{D}_k u_k \le\bar{g}, \quad  k=0,\dots,N-1,\\
& \underline h\le h_k(x_k,u_k)\le\bar h, \quad  k=0,\dots,N-1,\\
\end{aligned}
\end{equation}

where $x_{k+1}=f_k(x_k,u_k)$ is one integration step of the continuous dynamics over $[t_k,t_{k+1}]$ for notational simplicity,
$r(x_0,x_N) = 0$ represents the initial and terminal equality constraints,
$\mathcal{Q},\mathcal{R},\mathcal{Q}_N\succeq 0$ are weighting matrices, 
$\mathcal{S}_k$ is the state--control cost-coupling matrix,
and $x_k^{\mathrm{ref}},u_k^{\mathrm{ref}}$ are the reference points. 
$\mathcal{C}x_k + \mathcal{D}u_k$ includes the affine and box path constraints,
$h_k$ collects the nonlinear path constraints.

When using the exact Hessian of the Lagrangian in the SQP iterations, the resulting Hessian of the QP subproblem at stage $k = 0, \dots, N-1$ reads:
\begin{align}
    H_k &= \nabla^2 \mathcal{L}(x, u, \pi, \lambda)  \notag \\
        &= \begin{bmatrix} \mathcal{Q}_k & \mathcal{S}_k^{\top} \\ \mathcal{S}_k & \mathcal{R}_k \end{bmatrix} + \sum_{j=1}^{n_x} \pi_k \nabla^2 f_k(x_k,u_k) + \sum_{j=1}^{n_c} \lambda_k \nabla^2 h_k(x_k,u_k),
\end{align}
where $n_x$ is the dimension of states and $n_c$ is the number of nonlinear constraints, 
$\pi_k$ and $\lambda_k$ are the Lagrange multipliers associated with the dynamics and nonlinear constraints, respectively.

The stage cost is quadratic, so its Hessian $\left[\begin{smallmatrix}\mathcal{Q}_k & \mathcal{S}_k^{\top}\\ \mathcal{S}_k & \mathcal{R}_k\end{smallmatrix}\right]$ is constant and exact. 
The indefiniteness of $H_k$ thus has two sources: a constant contribution from the cost block, 
and an iterate-dependent contribution from the two curvature terms, 
which is usually the dominant one during the SQP iterations.
In these curvature terms, each component Hessian $\nabla^2[f_k]_j$ and $\nabla^2[h_k]_j$ is symmetric but carries no definiteness guarantee, 
and the associated multipliers $(\pi_k)_j$ and $(\lambda_k)_j$ may take arbitrary sign and magnitude during the SQP iterations. 
Their weighted sum is thus a general symmetric and typically indefinite matrix.
Adding it to the constant cost block can introduce negative eigenvalues in $H_k$, 
in particular when the multipliers are large, 
i.e.\ when constraints are strongly active or the current iterate is still far from a solution. 
Consequently, the exact-Hessian QP subproblem is in general nonconvex.

Importantly, this indefiniteness is confined to the individual stage blocks: 
$H_k$ remains block-diagonal across stages and, together with the banded Jacobian induced by the shooting constraints, preserves the
characteristic sparse structure of the OCP.

Collecting the stage Hessians $H_k$ and Jacobians and linearizing the constraints of~\eqref{eq:ocp} at the current SQP iterate yields a
structured QP. The subproblem takes the form:
\begin{equation}
\begin{aligned}
  & \min_{\substack{\Delta x_0,\dots,\Delta x_N\\ \Delta u_0,\dots,\Delta u_{N-1}}}
  && \sum_{k=0}^{N-1} \left( \tfrac{1}{2}
     \begin{bmatrix} \Delta x_k \\ \Delta u_k \end{bmatrix}^{\!\top}
     \begin{bmatrix} Q_k & S_k^\top \\ S_k & R_k \end{bmatrix}
     \begin{bmatrix} \Delta x_k \\ \Delta u_k \end{bmatrix}
     + \begin{bmatrix} q_k \\ r_k \end{bmatrix}^{\!\top}
     \begin{bmatrix} \Delta x_k \\ \Delta u_k \end{bmatrix} \right) \\
  & && +\ \tfrac{1}{2}\, \Delta x_N^\top Q_N\, \Delta x_N + q_N^\top \Delta x_N \\
  & \text{subject to}
  && \Delta x_0 = \Delta \bar{x}_0, \\
  & && \Delta x_N = \Delta \bar{x}_N, \\
  & && \Delta x_{k+1} = A_k \Delta x_k + B_k \Delta u_k + c_k, \quad k = 0, \dots, N-1, \\
  & && \begin{bmatrix} \underline{u}_k \\ \underline{x}_k \\ \underline{d}_k \end{bmatrix}
       \le
       \begin{bmatrix} J_k^{b,u} & 0 \\ 0 & J_k^{b,x} \\ D_k & C_k \end{bmatrix}
       \begin{bmatrix} \Delta u_k \\ \Delta x_k \end{bmatrix}
       \le
       \begin{bmatrix} \overline{u}_k \\ \overline{x}_k \\ \overline{d}_k \end{bmatrix},
       \quad k = 0, \dots, N-1
\end{aligned}
\end{equation}
where, at each stage $k$, 
$\Delta x_k \in \mathbb{R}^{n_x}$ and $\Delta u_k \in \mathbb{R}^{n_u}$
denote the state and control deviation variables, and $N$ is the horizon length.

All remaining quantities are obtained by linearizing the OCP at the
current SQP iterate $(x_k,u_k)$ with multipliers $(\pi_k,\lambda_k)$:
\begin{itemize}
    \item $Q_k$, $R_k$, and $S_k$ are the state, control, and state--control coupling blocks of the exact Lagrangian Hessian $H_k$.
    $q_k=\mathcal{Q}_k(x_k-x_k^{\mathrm{ref}})+\mathcal{S}_k^{\top}(u_k-u_k^{\mathrm{ref}})$ and 
    $r_k=\mathcal{R}_k(u_k-u_k^{\mathrm{ref}})+\mathcal{S}_k(x_k-x_k^{\mathrm{ref}})$ are the corresponding stage-cost gradients, 
    and $Q_N$, $q_N$ are the terminal Hessian and gradient.

    \item $A_k=\nabla_{x_k} f_k$ and $B_k=\nabla_{u_k} f_k$ are the
    Jacobians of the discretized dynamics, and $c_k=f_k(x_k,u_k)-x_{k+1}$ is the dynamics (shooting) residual at the current iterate; 
    together they define the linearized state-transition equation. At a dynamically feasible iterate $c_k=0$.

    \item $\Delta\bar{x}_0=\bar{x}_0-x_0$ and $\Delta\bar{x}_N=\bar{x}_N-x_N$ are the residuals of the initial and terminal equality constraints, 
    i.e.\ the corrections needed to satisfy $x_0=\bar{x}_0$ and $x_N=\bar{x}_N$ after the step. 
    The terminal equality is present only for the problems that impose it; otherwise this row is absent.

    \item $J_k^{b,u}$ and $J_k^{b,x}$ are selection matrices picking out the bounded control and state components, 
    with bounds $\underline{u}_k,\overline{u}_k$ and $\underline{x}_k,\overline{x}_k$ (box constraints).

    \item $C_k=\nabla_{x_k} h_k$ and $D_k=\nabla_{u_k} h_k$ are the Jacobians of the general path constraints, 
    with $\underline{d}_k=\underline{h}-h_k(x_k,u_k)$ and $\overline{d}_k=\overline{h}-h_k(x_k,u_k)$; 
    these constraints couple states and controls. For the affine constraints $\mathcal{C}x_k+\mathcal{D}u_k$ of~\eqref{eq:ocp}, 
    $C_k$ and $D_k$ reduce to the constant matrices $\mathcal{C}$ and $\mathcal{D}$.
\end{itemize}

\subsection{Benchmark OCPs}
In this subsection, we introduce the formulation of the eight OCPs used to construct the benchmark. 
Each OCP is obtained by taking a dynamic system provided in~\cite{Kiessling2026} 
and casting it as a regulation or point-to-point transfer problems with problem-specific modifications,
which can be described as \eqref{eq:ocp}.
The eight OCPs are named and characterized as follows:
\begin{itemize}
    \item \textbf{Pendulum:} An inverted pendulum on an actuated cart,
    driven by a horizontal force. It is regulated from the downward to the
    upright equilibrium subject to an input bound and an additional mixed
    state--input inequality constraint. Its stage cost further includes a
    bilinear state--input term, which facilitates the occurrence of indefinite Hessians in the SQP iterations.
    \item \textbf{Acrobot:} Adapted from \cite{spong1995swing}. 
    A two-link pendulum actuated only at the elbow
    joint, with joint friction. The state is driven between two prescribed
    configurations through terminal equality constraints, yielding a strongly
    nonlinear, underactuated maneuver.
    \item \textbf{Dubins car:} Adapted from \cite{dubins1957curves}.
    A kinematic unicycle with position and heading as
    states, actuated by forward velocity and turn rate. It performs a
    parallel-parking maneuver, a lateral displacement returning to the initial
    heading, within a narrow corridor imposed by state-box constraints, with a
    terminal equality constraint on the full pose.
    \item \textbf{Furuta pendulum:} Adapted from \cite{furuta1992swing}.
    A rotary inverted pendulum with
    experimentally identified parameters and joint damping, actuated by a single
    base-arm torque. The pendulum is swung up and stabilized at the upright
    position under tight torque limits and a non-uniform time grid.
    \item \textbf{Flying robot:} Adapted from \cite{betts2010practical}.
    A planar free-flying robot with position,
    orientation and their rates, propelled by four non-negative thrusters that
    together produce a body-aligned net force and a turning moment. The state is
    driven toward a reference point---zero position and orientation with a
    prescribed nonzero translational velocity---under thrust non-negativity and
    summed-thrust linear constraints.
    \item \textbf{Robot arm:} Adapted from \cite{betts2010practical}.
    A three-DOF robot arm (radial extension, rotation,
    elevation) with configuration-dependent inertia and a cross-axis actuator
    coupling that is nonlinear in the inputs. A joint angle is regulated to a new
    set-point under state and input bounds; the input coupling renders the exact
    Hessian indefinite.
    \item \textbf{Crane:} Adapted from \cite{betts2010practical}.
    An overhead crane with hoisting length, trolley
    position and payload swing angle, actuated by trolley and hoisting
    accelerations. The load is transferred between two rest positions while
    respecting a swing-rate bound and a nonlinear region constraint on the
    payload.
    \item \textbf{Moonlander:} Adapted from \cite{meditch1964optimal}.
    A planar lunar lander with translational and
    rotational motion under reduced gravity, actuated by two laterally offset
    thrusters. It performs a soft landing, reaching a ground target with zero
    terminal velocity, under non-negative, bounded thrust.
\end{itemize}

\begin{table}[H]
    \centering
    \caption{Summary of the evaluated OCP problems. The state dimension
    $n_x$, control dimension $n_u$, and horizon length $N$ span a range
    of common control problem scales.}
    \label{tab:ocp_problems}
    \begin{tabular}{l ccc}
        \toprule
        \textbf{Problem} & $n_x$ & $n_u$ & $N$ \\
        \midrule
        Pendulum     & 4 & 1 & 25 \\
        Acrobot      & 4 & 1 & 15 \\
        Dubins car   & 3 & 2 & 25 \\
        Furuta       & 4 & 1 & 20 \\
        Flying robot & 6 & 4 & 50 \\
        Robot arm    & 6 & 3 & 30 \\
        Crane        & 6 & 2 & 40 \\
        Moonlander   & 6 & 2 & 30 \\
        \bottomrule
    \end{tabular}
\end{table}

Besides the different physical systems,
These OCPs also cover different state and control dimensions, as shown in
Table \ref{tab:ocp_problems}. The state dimension $n_x$ ranges from $3$ to
$6$, the control dimension $n_u$ from $1$ to $4$, and the horizon length $N$ from $15$ to $50$.
This spread is chosen so that the benchmark spans a range of common problem scales in control problems,
which allows the evaluated methods to be assessed across problems of varying size and structure.

\subsection{Collecting Strategy}
Based on the formulated OCPs, 
we can construct a representative benchmark of indefinite QPs arising from SQP-based NMPC,
avoiding the artificial assumptions of randomly generated Hessians and constraints. 
Specifically, we propose a collection strategy that extracts the QP subproblems directly from intermediate SQP iterations. 
By logging the linearized dynamics, stage constraints, and indefinite objective approximations exactly as they arise in each iteration, 
this procedure ensures that every benchmark instance faithfully reflects the numerical conditioning and sparsity encountered in practice.

The collection strategy consists of two steps: \emph{generate} and \emph{collect}, and can be summarized as follows:

\begin{enumerate}
    \item \textbf{Generate:} In this step, we generate iterates by solving
    the OCP. To obtain iterates at various points during the solution
    process, we can apply SQP-level regularization techniques (such as PROJECT or MIRROR) or utilize the Gauss-Newton Hessian approximation in \texttt{acados}.

    \item \textbf{Collect:} In this step, we utilize the previously generated iterates as initial guesses for linearization within the solver. 
    By restricting the maximum number of SQP iterations, we can selectively collect indefinite QPs that exhibit an indefinite reduced Hessian.
\end{enumerate}

In the collect step, QPs are all with two failling status code from \texttt{acados}, code \texttt{2} for max iterations reached and code \texttt{4} for infeasible QP.

The QPs with code \texttt{2} are from limiting the number of SQP iterations to prevent the solver moving towards the neighborhood of the OCP solution,
the function \texttt{qp\_diagnostics} from the class \texttt{AcadosOcpSolver} is called to check the definiteness of the reduced Hessian, 
and only the QPs with indefinite reduced Hessians are collected.

The QPs with code \texttt{4} are from the failing of the Riccati recursion in \texttt{HPIPM},
this is achieved by setting the Riccati algorithm in HPIPM as classical Riccati recursion, i.e. \texttt{ric\_alg = 0}, 
instead of the square-root version \cite{frisonAlgorithmsMethodsHighPerformance},
which requires the reduced Hessian to be positive definite in order to proceed with the Riccati recursion.
So when the HPIPM fails to proceed with the Riccati recursion, it indicates that the reduced Hessian is indefinite, and the corresponding QP is collected.

For each problem, we limit the number of near-identical QPs, 
where similarity is measured by the minimal eigenvalue and the rate of indefinite stages (both introduced in Section~\ref{sec:benchmark_analysis}), 
so that the number of instances per problem stays within a comparable range and the collected instances remain meaningful.
We further use IPOPT~\cite{wachter2006implementation} as a reference solver and keep only the QPs that it solves successfully. 
After this filtering, we obtain a total of $175$ indefinite QPs from the eight OCPs,
the per-problem counts are reported in Table~\ref{tab:collected_qps} and range from $15$ to $29$.

\begin{table}[H]
    \centering
    \caption{Number of indefinite QPs collected from each OCP ($175$ in total).}
    \label{tab:collected_qps}
    \begin{tabular}{l c}
        \toprule
        \textbf{Problem} & \textbf{Number} \\
        \midrule
        Pendulum     & 18 \\
        Acrobot      & 15 \\
        Dubins car   & 29 \\
        Furuta       & 19 \\
        Flying robot & 20 \\
        Robot arm    & 24 \\
        Crane        & 29 \\
        Moonlander   & 21 \\
        \bottomrule
    \end{tabular}
\end{table}

The QPs are stored in \texttt{json} format and are available at \url{https://github.com/acados/ocp_qp_dataset_collection}.

\section{Benchmark Analysis}
\label{sec:benchmark_analysis}
In this section,
we will introduce some analysis of the collected problem instances
for gaining a better understanding of characteristics regarding definiteness in the indefinite OCP QP benchmark,
The distribution of indefiniteness is analyzed in Section \ref{subsec:distribution},
which describes the magnitude and coverage of the indefinite reduced Hessians in the collected QPs,
and the source of indefiniteness is analyzed in Section \ref{subsec:source}, 
which further illustrates the source of indefiniteness in the collected QPs by analyzing the eigenvalues of the stage Hessians.

\subsection{Distribution of Indefiniteness}
\label{subsec:distribution}
We first characterize the collected QPs along two axes:
the \emph{depth} of indefiniteness, measured by the magnitude of the smallest (most negative) eigenvalue $|\lambda_{\min}|$ of the reduced Hessian,
and the \emph{width}, measured by the fraction of block reduced Hessians that are indefinite within a QP (the indefinite-stage rate).

\begin{figure}[H]
    \centering
    \includegraphics[width=0.95\textwidth]{images/benchmark/indefiniteness_coverage_scatter.pdf}
    \caption{Depth ($|\lambda_{\min}|$ of the reduced Hessian, log scale)
    versus width (indefinite-stage rate) of the collected QPs, coloured by
    OCP problem.}
    \label{fig:indefinite_qp_scatter}
\end{figure}

As seen from Figure~\ref{fig:indefinite_qp_scatter}, the collected indefinite QPs do not concentrate in a single regime.
They spread over more than seven orders of magnitude in depth ($|\lambda_{\min}|$ ranging from about $10^{-3}$ to $10^{4}$).
And the coverage is strongly problem-dependent: each OCP occupies a characteristic region of the depth--width plane rather than scattering uniformly.

Two complementary major patterns can be observed.
Problems that are indefinite in only a small fraction of their stages, such as \texttt{overhead\_crane},
\texttt{free\_flying\_robot} and \texttt{furuta\_cons}, nevertheless reach very deep negative curvature:
their indefinite-stage rates stay below $0.30$ (with medians of $0.05$, $0.02$ and $0.10$, respectively),
while $|\lambda_{\min}|$ reaches up to $10^{2}$--$10^{4}$.
Conversely, problems with a high indefinite-stage rate, \texttt{dubins\_car\_parallel\_parking} and \texttt{robot},
with rates ranging from $0.47$--$1.0$, tend to exhibit only moderate depth,
with $|\lambda_{\min}|$ clustered around $10^{-1}$--$10^{1}$.
Broad and deep indefiniteness thus appear to occur largely separately rather than together.

Some problems also fall in the middle, such as \texttt{moonlander} and
\texttt{acrobot}, which occupy intermediate indefinite-stage rates
(roughly $0.07$--$0.53$). \texttt{moonlander} keeps its depth in a narrow band
($|\lambda_{\min}|$ between $10^{0}$ and $10^{2}$), whereas \texttt{acrobot}
scatters over almost the entire depth range, from near-singular values around
$10^{-2}$ up to nearly $10^{4}$. They neither confine their negative curvature
to a few stages nor spread it across the whole horizon, and thus bridge the two
extreme regimes rather than belonging clearly to either.

Taken together, the eight OCPs cover a broad range of the depth--width plane:
from narrow-but-deep (\texttt{overhead\_crane}, \texttt{free\_flying\_robot}, \texttt{furuta\_cons}),
through intermediate (\texttt{moonlander}, \texttt{acrobot}),
to broad-but-moderate (\texttt{dubins\_car\_parallel\_parking}, \texttt{robot}),
with \texttt{pendulum\_cons} adding cases of very shallow negative curvature
($|\lambda_{\min}|$ mostly below $10^{-1}$) at intermediate indefinite-stage rates.
This spread indicates that the collected QP subproblems exercise indefiniteness across qualitatively different regimes rather than a single one.

\begin{figure}[H]
    \centering
    \begin{minipage}[t]{0.48\textwidth}
        \centering
        \includegraphics[width=\linewidth]{images/benchmark/indefiniteness_coverage_depth.pdf}
        \subcaption{Depth: distribution of $|\lambda_{\min}|$.}
        \label{fig:indefinite_depth}
    \end{minipage}
    \hfill
    \begin{minipage}[t]{0.48\textwidth}
        \centering
        \includegraphics[width=\linewidth]{images/benchmark/indefiniteness_coverage_breadth.pdf}
        \subcaption{Width: distribution of the indefinite-stage rate.}
        \label{fig:indefinite_breadth}
    \end{minipage}
    \caption{Marginal distributions of depth and width over all collected QPs.}
    \label{fig:indefinite_marginals}
\end{figure}

Figure~\ref{fig:indefinite_marginals} shows the two marginal distributions of the same data. The depth histogram
(Figure~\ref{fig:indefinite_depth}) is centred around $|\lambda_{\min}| \approx 10$, where most of the indefinite QPs lie,
with roughly one third of all instances falling between $10^{0.5}$ and $10^{1.5}$,
and decays toward both a near-singular tail near $10^{-3}$ and a small number of very deep cases at $10^{3}$--$10^{4}$,
of which only five instances exceed $10^{3}$.
Hence, although extreme cases exist, the negative curvature typically encountered during the SQP iterations is of moderate magnitude.

The width histogram (Figure~\ref{fig:indefinite_breadth}) shows how the indefinite-stage rate is distributed over the collected QPs. The rate
covers the full range up to~$1$, with a large share of QPs being indefinite in only a minority of their stages:
close to $40\%$ of the instances have a rate below $0.11$. The remainder is spread out
with a secondary accumulation around rates of $0.33$--$0.55$ and a further group around $0.72$--$0.88$,
while only a single QP is indefinite in essentially all of its stages.

This is consistent with the scatter plot: for most problems only a subset of stages become indefinite, whereas a few problems
(\texttt{dubins\_car\_parallel\_parking}, \texttt{robot}) are indefinite across the majority of their stages.

\subsection{Source of Indefiniteness}
\label{subsec:source}
Recalling the block reduced Hessian at $k$ stage is calculated as $M_k = R_k + B^\top P_{k+1}B$,
which is assembled from two contributions: the control cost block $R$ and the cost-to-go block $P$. 

To understand \emph{where} the indefiniteness originates,
we classify every indefinite stage according to which of the two blocks has a negative smallest eigenvalue: 
\emph{R only} ($R$ indefinite, $P$ definite), \emph{P only} ($P$ indefinite, $R$ definite), and \emph{both} (both indefinite).

\begin{figure}[H]
    \centering
    \includegraphics[width=0.7\textwidth]{images/benchmark/stage_indefiniteness_source.pdf}
    \caption{Number of indefinite stages by source: control cost~$R$,
    cost-to-go~$P$, or both.}
    \label{fig:indefinite_source_counts}
\end{figure}

Figure~\ref{fig:indefinite_source_counts} aggregates all indefinite stages across the benchmark. 
The dominant source is the joint case: $1114$ of the $1571$ indefinite stages ($71\%$) have \emph{both} $R$ and $P$ indefinite, 
so in the large majority of cases the negative curvature is not confined to a single block. 
Among the remaining stages where only one block is indefinite, 
the cost-to-go block is the more frequent culprit: $P$ alone accounts for $299$ stages ($19\%$), 
nearly twice the $158$ stages ($10\%$) in which only the control cost $R$ is indefinite. 

Overall, the cost-to-go block $P$ is indefinite more often than the control cost block $R$,
being involved in $1413$ of the $1571$ indefinite stages ($90\%$) against $1272$ ($81\%$) for $R$,
which indicates that the indefiniteness is driven primarily by the cost-to-go term rather than by the control cost.

\begin{figure}[H]
    \centering
    \includegraphics[width=1.0\textwidth]{images/benchmark/stage_indefiniteness_signature.pdf}
    \caption{Per-problem composition of stage indefiniteness. Each bar
    shows, for one OCP, the share of stages that are \emph{R only},
    \emph{both}, \emph{P only}, ($n$ is the number of collected indefinite stages).}
    \label{fig:indefinite_source_signature}
\end{figure}

While Figure~\ref{fig:indefinite_source_counts} gives the aggregate picture,
Figure~\ref{fig:indefinite_source_signature} breaks the same classification down per problem and reveals a strongly problem-dependent and clean structure: 
six of the eight OCPs are dominated by a single category ($\geq\!70\%$ of their stages),
three of them almost exclusively so ($\geq\!90\%$),
and only \texttt{acrobot} and \texttt{overhead\_crane} are genuinely mixed.

As observed from Figure~\ref{fig:indefinite_source_signature},
most of the problems are governed by the \emph{both} case, 
including \texttt{moonlander} ($96\%$), \texttt{robot} ($92\%$), \texttt{dubins\_car\_parallel\_parking} ($76\%$)
and \texttt{furuta\_cons} ($72\%$).
Notably, three problems (\texttt{dubins\_car\_parallel\_parking}($n=581$), \texttt{robot}($n=404$) and \texttt{moonlander}($n=198$)), contribute by far the most indefinite stages, 
they largely account for the aggregate dominance of the \emph{both} case reported above. 

The \emph{P only} case dominates the remaining problems: \texttt{free\_flying\_robot} is indefinite purely through the cost-to-go block ($100\%$),
and \texttt{pendulum\_cons} nearly so ($87\%$). The two mixed problems lean the same way:
\texttt{overhead\_crane} splits into $58\%$ \emph{P only} against $39\%$ \emph{both},
and \texttt{acrobot} into $54\%$ \emph{P only} against $41\%$ \emph{both},
so that even where no single category dominates, the cost-to-go block is the more frequent one.
Taken over all stages, the cost-to-go block is indefinite in the large majority of cases, 
whereas the \emph{R only} case is marginal:
it appears in only four problems and is negligible for \texttt{robot}, \texttt{acrobot} and \texttt{overhead\_crane} (all below $5\%$), 
with \texttt{dubins\_car\_parallel\_parking} being the sole problem where isolated $R$-block indefiniteness is substantial ($24\%$ of its stages,
accounting for $140$ of the $158$ \emph{R only} stages in the whole benchmark). 
Indefiniteness therefore predominantly involves the cost-to-go block, 
either alone or jointly with $R$, while $R$-block-only indefiniteness is the exception.

In conclusion, these figures numerically shows that,
across the benchmark the cost-to-go block is involved in the vast majority of indefinite stages, 
a regularization strategy that merely enforces $M\succ 0$ would therefore leave the negative curvature unaddressed in most stages. 
Since the cost-to-go term enters the reduced Hessian through $M_k = R_k + B^\top P_{k+1} B$, 
keeping $P_{k+1}$ (and hence $M_k$) definite is essential for a reliable regularization method, 
whereas treating the control cost block alone is not sufficient.

\end{document}